\section{Experiments}
\label{sec:experiments}

We evaluate \sysname{} along five research questions:
\begin{itemize}[leftmargin=*]
\item \textbf{RQ1:} Does \sysname{} improve severity-label agreement over strong baselines?
\item \textbf{RQ2:} Do blinded specialists rate \sysname{} reports more highly, especially on evidence sufficiency?
\item \textbf{RQ3:} How much do the temporal gate and agentic loop each contribute?
\item \textbf{RQ4:} Does validity gating improve Product--Litigation Risk Entailment?
\item \textbf{RQ5:} Does the calibrated stopping rule control evidence-omission risk under temporal shift?
\end{itemize}

\subsection{Experimental Setup}
\label{sec:setup}

\paragraph{Dataset split and label provenance.}
VIGIL-Dataset is partitioned 70/10/20 into train/validation/test splits stratified by year and litigation type (not by dispute identity; Limitations).
Because a substantial share of labels originate from an automated pre-labeler and only a subset are human-adjudicated, primary claims are computed exclusively on the human-adjudicated test split ($n = 82$ cases); retained-label results are reported solely for scale.
Conformal calibration draws on human-adjudicated cases from the training and validation pools, kept disjoint from the test cases.
For \sysname{} and baselines B2--B5 the training split is used only to select few-shot examples and tune thresholds; no LLM is fine-tuned.

\plre{} pairs (1{,}856 total; 28.4\% positive; 19.7\% stale-authority distractors) are partitioned 1{,}112 / 372 / 372 (train / validation / test; $\approx$60/20/20), stratified by product archetype and year so that every archetype appears in every split.
Few-shot exemplars and the entailment threshold use only train/validation; all reported metrics use the 372 held-out test pairs.

\paragraph{Baselines.}
We compare \sysname{} against five risk-classification baselines: B1 Additive-only Rule (the pre-labeler's scoring matrix without its composite bonus, thresholds tuned on validation); B2 Single-Pass LLM; B3 RAG-only; B4 Single-Agent (\sysname{} without the goal-conditioned loop); and B5 Iteration-Matched Self-Refinement (Reflexion-style~\citep{shinn2023reflexion}, matched on \sysname{}'s mean budget of 2.4 iterations, not on tokens).
B2--B5 share the same GPT-4o backbone and decoding settings as \sysname{}.
For \plre{} we further include non-VIGIL baselines that isolate suspected confounds: product--theory cosine, LegalBERT product--theory NLI, COLIEE-style~\citep{rabelo2022coliee} precedent retrieval without interval closure, and \sysname{} without temporal gating.

\paragraph{Automatic evaluation.}
All reported automatic metrics freeze the treatment graph produced by the extractor; low-confidence edges are not human-reviewed at evaluation time.
We report macro F1/precision/recall over three severity labels, report-quality ratings by two blinded specialists, loop efficiency, \plre{} precision/recall/F1, edge-error propagation, and evidence-coverage omission risk under calibrated thresholds.
Severity-label F1 measures triage-label agreement, not legal liability.

\subsection{RQ1: Severity-Label Agreement}
\label{sec:rq1}

Table~\ref{tab:severity} reports severity-label agreement on the human-adjudicated subset ($n = 82$).
\sysname{} reaches 0.79 macro F1 vs.\ 0.75 for iteration-matched B5; the same rank order holds on the full held-out test set.
With $n = 82$ we do not treat this four-point gap as established; these triage numbers are secondary to \plre{} and are not significance-tested.

\begin{table}[t]
\centering
\small
\setlength{\tabcolsep}{4pt}
\begin{adjustbox}{max width=\columnwidth}
\begin{tabular}{lcccc}
\toprule
System & \multicolumn{3}{c}{Human-adj.} & Full \\
\cmidrule(lr){2-4} \cmidrule(lr){5-5}
 & F1 & P & R & F1 \\
\midrule
B1 Additive-only & 0.55 & 0.55 & 0.56 & 0.62 \\
B2 Single-Pass LLM & 0.70 & 0.70 & 0.71 & 0.72 \\
B3 RAG-only & 0.72 & 0.72 & 0.72 & 0.74 \\
B4 Single-Agent & 0.74 & 0.74 & 0.74 & 0.76 \\
B5 Self-Refinement & 0.75 & 0.75 & 0.75 & 0.77 \\
VIGIL Full & \textbf{0.79} & \textbf{0.79} & \textbf{0.80} & \textbf{0.81} \\
\bottomrule
\end{tabular}
\end{adjustbox}
\caption{Macro severity-label agreement.
Human-adjudicated subset ($n = 82$; primary).
Full: held-out test including retained pre-labels.}
\label{tab:severity}
\end{table}

\subsection{RQ2: Report Quality}
\label{sec:rq2}

Two blinded specialists rated \sysname{} and B2--B5 reports on 15 cases (75 reports) along Accuracy, Completeness, Actionability, Evidence Sufficiency, and Clarity (5-point scales).
\sysname{}'s mean scores are 4.3/4.1/4.2/4.4/3.9 on the five scales, highest on Evidence Sufficiency.
Mean inter-rater $\kappa = 0.68$ (0.62--0.74 across scales).
RQ2 is supporting evidence only ($n = 15$): we make no pairwise superiority claim over B2--B5, and any Wilcoxon signed-rank comparison on per-case means is exploratory and uncorrected for multiplicity.

\subsection{RQ3: Agentic Loop and Temporal Graph Ablations}
\label{sec:rq3}

Temporal precedent reasoning is the central empirical claim.
Table~\ref{tab:ablation} isolates the two mechanisms of Eq.~\eqref{eq:gate}.
Full validity-interval gating plus treatment weighting yields \plre{} F1 of 0.78 and a stale-citation rate of 4.1\%.
Removing treatment weighting alone drops F1 by 0.03 and roughly doubles stale citations (7.8\%).
Removing interval gating entirely (recovering flat semantic similarity) widens the gap to $\Delta$F1 $= -0.06$ and raises the stale-citation rate to 12.7\%---a three-fold increase.
Without the gate, retrieval admits more lexically similar but already-closed authorities, altering the evidence that reaches the generator.
This controlled degradation constitutes the paper's main methodological result.

\begin{table}[t]
\centering
\small
\setlength{\tabcolsep}{4pt}
\begin{adjustbox}{max width=\columnwidth}
\begin{tabular}{lccc}
\toprule
Variant & F1 & $\Delta$F1 & Stale \% \\
\midrule
VIGIL Full (gate + weight) & 0.78 & --- & 4.1 \\
$-$ treatment weighting $w(\cdot)$ & 0.75 & $-0.03$ & 7.8 \\
$-$ interval gating (flat sim.) & 0.72 & $-0.06$ & 12.7 \\
\bottomrule
\end{tabular}
\end{adjustbox}
\caption{Temporal precedent-graph ablation (main claim).
Both mechanisms of Eq.~\eqref{eq:gate} are removed in sequence; the bottom row recovers flat semantic similarity.
F1 is \plre{} F1; Stale is the stale-citation rate, i.e., citations to already-closed precedent ($t_q \ge e_v$).}
\label{tab:ablation}
\end{table}

Negative-treatment error propagation was measured by substituting gold labels on the 600 annotated spans and by additionally masking extracted edges (Appendix~\ref{app:propagation}).
The extractor operates at the observed miss rate ($\approx$23\%, from 0.77 recall) and costs 2 \plre{} F1 points relative to gold-span substitution (0.80, stale 3.7\%).
Masking a further 30\% of extracted negative-treatment edges ($\approx$46\% effective miss rate) returns F1 to the no-gate level (0.72) with 11.3\% stale citations, so the gate's benefit depends on treatment-edge quality.
Stale citations persist even with gold edges because a closed authority can enter context through precedent chaining (\S\ref{sec:rag}) rather than as gated support, and because annotator validity judgments can disagree with interval closure on boundary cases.

\subsection{RQ4: Product--Litigation Risk Entailment}
\label{sec:rq4}

\plre{} uses anonymized product-attribute cards (five archetypes) paired with docket-grounded theory cards across 312 product-mapping cases (1{,}856 pairs; 28.4\% positive; 19.7\% stale-authority distractors; $\kappa = 0.71$).
Table~\ref{tab:plre}: VIGIL Full reaches 0.78 F1 ($+$0.06 vs.\ no-gate; $+$0.05--0.13 versus non-VIGIL product--theory and legal-NLI baselines).
VIGIL without gating (0.72) sits slightly below RAG-only (0.73), suggesting that the agentic loop without a validity filter can reinforce stale but lexically close precedent; in this ablation, the \plre{} gain comes with the gate.
Gains partly reflect rejection of stale distractors (19.7\% of pairs); this version reports F1 on the full held-out set only, not on the non-distractor subset (Limitations).
Primary claims rest on the 372 held-out pairs.

\begin{table}[t]
\centering
\small
\setlength{\tabcolsep}{4pt}
\begin{adjustbox}{max width=\columnwidth}
\begin{tabular}{lccc}
\toprule
System & P & R & F1 \\
\midrule
Product--theory cosine & 0.64 & 0.66 & 0.65 \\
LegalBERT NLI & 0.69 & 0.68 & 0.68 \\
COLIEE-style, no gate & 0.71 & 0.70 & 0.70 \\
RAG-only generation & 0.73 & 0.74 & 0.73 \\
VIGIL w/o gating & 0.72 & 0.73 & 0.72 \\
VIGIL Full & \textbf{0.77} & \textbf{0.80} & \textbf{0.78} \\
\bottomrule
\end{tabular}
\end{adjustbox}
\caption{\plre{} comparison against external legal-NLI and precedent-retrieval baselines (LegalBERT NLI: product--theory NLI).
Held-out test pairs only.}
\label{tab:plre}
\end{table}

\subsection{RQ5: Evidence-Coverage Risk Control and Temporal Shift}
\label{sec:rq5}

Calibrating the evidence threshold to $\alpha \in \{0.05, 0.10, 0.20\}$ yields omission risk 0.043/0.082/0.171 in random splits.
Under a temporal split (2023--24 $\to$ 2025) the achieved risk rises to 0.061/0.113/0.196, exceeding the target at $\alpha = 0.05$ and 0.10, which motivates rolling re-calibration (Appendix~\ref{app:theory}).

\paragraph{Backbone dependence.}
Replacing GPT-4o with an open-weights Llama-3.1-70B~\citep{grattafiori2024llama3} + bge~\citep{xiao2024cpack} stack, unchanged otherwise, lowers full-test-set severity macro F1 from 0.81 to 0.76 (mean over five runs) with 2.6 vs.\ 2.4 mean iterations (Appendix~\ref{app:theory}).
Baselines were not re-run on this stack, so this shows that \sysname{} runs without the proprietary model, not that its advantage is backbone-independent.
